\documentclass[a4paper,11pt]{amsart}
\usepackage[british]{babel}
\usepackage[colorlinks=true,citecolor=black,linkcolor=black]{hyperref}
\usepackage{cite,tikz,amssymb}
\usepackage[notcite, notref, color]{showkeys}
\usepackage[normalem]{ulem}

%%%%%%% Draft watermark
\usepackage[ddmmyy]{datetime}
\usepackage{tikz,tikz-cd, tikzpagenodes, lmodern}
\usepackage[pages=all]{background}
\backgroundsetup{angle=0, opacity=0.5, scale=1, contents= {   
\begin{tikzpicture}[remember picture, overlay, scale=3]           
\fontsize{20}{10}\selectfont
\node[text=red,rotate=270, above=1.5cm] at (current page text area.east) 
{\textsc{Draft version -- {\today}}};
\end{tikzpicture} }}
%%%%%%%%%%%%


%%%%%%%%%% Theorems
\newtheorem{theorem}{Theorem}
\newtheorem{proposition}[theorem]{Proposition}
\newtheorem{corollary}[theorem]{Corollary}
\newtheorem{lemma}[theorem]{Lemma}
\newtheorem{definition}[theorem]{Definition}
\theoremstyle{remark}
\newtheorem{remark}[theorem]{Remark}

%%%%%%%%%% Custom notation
\newcommand{\abs}[1]{\left|#1\right|} % Absolute value
\newcommand{\norm}[1]{\left\|#1\right\|} % Norm
\newcommand{\id}{\operatorname{I}} % The identity
\newcommand{\cB}{\mathcal{B}} % Banach space
\newcommand{\cC}{\mathcal{C}} % Banach space
\newcommand{\bC}{\mathbb{C}} % Complex numbers
\newcommand{\bN}{\mathbb{N}} % Natural numbers
\newcommand{\bR}{\mathbb{R}} % Reals
\newcommand{\bZ}{\mathbb{Z}} % Integers
\newcommand{\Ker}{\operatorname{Ker}} % Kernel
\newcommand{\Ran}{\operatorname{Ran}} % Range
\newcommand{\Coker}{\operatorname{Coker}} % Cokernel
\newcommand{\EssSpec}{\operatorname{EssSpec}} % Essential spectrum
\newcommand{\nc}[1]{{\left|#1\right|}_{\lozenge}} % Ball measure of non-compactness
\newcommand{\EssRad}{\operatorname{EssRad}} % Essential spectral radius
\newcommand{\cG}{\mathcal{G}} % 
\newcommand{\cL}{\mathcal{L}} % 
\newcommand{\cR}{\mathcal{R}} % 


\begin{document}
\title{Measure of non-compactness and the essential spectral radius}
%\date{}

\author{Oliver Butterley}
\address{(Oliver Butterley) Dipartimento di Matematica, II Università di Roma (Tor Vergata) -- Via della Ricerca Scientifica -- Roma -- 00133 -- Italy}
\email{butterley@mat.uniroma2.it}


\begin{abstract}
    An estimate on the measure of non-compactness implies and estimate on the essential spectral radius.
\end{abstract}

\maketitle
% \thispagestyle{empty}

%%%%%%%%%%%%%%%%%%%%%%%%%%%%%%%%%%%%%%%%%%%%%%%%%%%%%%%%%%%%%%%%%%%%%%%%%%%%%



A bounded operator \(A : \cB \to \cC\) between two Banach spaces is said to be a \emph{Fredholm} operator if its kernel \(\Ker A\) and cokernel \(\Coker A := \cC / \Ran A\) are both finite-dimensional.
We say that \(\lambda\) lies in the essential spectrum \(\EssSpec(A)\) of a bounded operator \(A\) if \((\lambda \id - A)\) is not a Fredholm operator. 
The \emph{ball measure of non-compactness} of a bounded operator \(A: \cB \to \cC\) is denoted \(\nc{A}\) and is defined as the infimum of real numbers \(\epsilon>0\) such that the image of the unit ball under \(A\) admits a finite cover by open balls of radius \(\epsilon\).
% (\(\left\{ x + y : x \in \cB, \norm{x} \leq \epsilon \right\}\) for some \(y\in \cB\)).
The quantity \(\nc{\cdot}\) is a seminorm on the space of operators, is submultiplicative and \(\nc{A} = 0\) if and only if \(A\) is compact.

\begin{lemma}
    \label{lem:lim-exists}
    If \(A\) is a bounded operator on \(\cB\) then the limit 
    \(\rho := \lim_{n\to \infty} \nc{A^n}^{1/n}\)
    exists and satisfies \(0\leq \rho \leq \nc{A}\).
\end{lemma}
\begin{proof}
    Let \(p(n):= \log \nc{A^n}\) and observe that \(p\) is subadditive in the sense that \(p(m+n) \leq p(m)+p(n)\) for all \(m,n\in\bN\).
    The lemma then follows from \cite[Lemma 4.1.2]{Davies07}.
\end{proof}

\begin{lemma}
    \label{lem:dual-nc}
    If \(A\) is a bounded operator on \(\cB\) then \(\nc{A^*} \leq 2\nc{A}\).
\end{lemma}


\begin{lemma}
    \label{lem:is-fredholm}
    Let \(A\) be a bounded operator on \(\cB\).
    Then \((\lambda \id - A)\) is Fredholm whenever \(\abs{\lambda} > \displaystyle\lim_{n\to \infty}{\nc{A^n}}^{{1}/{n}}\).
\end{lemma}
\begin{proof}
    Replacing \(A\) with \(\lambda^{-1}A\) it suffices to prove that  \((\id - A)\) is Fredholm whenever \( \lim_{n}{\nc{A^n}}^{{1}/{n}} < 1\).

    Firstly we show that \uline{\(\cL := \Ker(\id -A )\) is finite dimensional}.
    Fix \(n\in \bN\) sufficiently large such that \(\lim_{n}{\nc{A^n}} < \frac{1}{4}\).
    In search of contradiction we suppose that \(\cL := \Ker(\id -A )\) fails to be finite dimensional.
    This means that there exists an infinite sequence \(x_j \in \cL\) such that \(\norm{x_j} = 1\) and \(\norm{x_j - x_k}\geq \frac{3}{4}\) whenever \(j\neq k\).
    Since \(A^n x_j = x_j\) for each \(j\), \(\norm{A^n x_j - A^n x_k} =  \norm{x_j - x_k} \geq \frac{3}{4}\) whenever \(j\neq k\).
    On the other hand, since \(\cup_{j}x_{j}\) is a subset of the unit ball and \(\nc{A^n} < \frac{1}{4}\), we know that \(\cup_{j} A^n x_j\) can be covered by a finite number of \(\frac{1}{4}\)-balls.
    This is a contradiction.
    
    Secondly we show that \uline{\(\cR:=\Ran(\id -A )\) is closed}.
    Let \(g_k \in \cR\) be a sequence which converges to \(g\).
    We need to show that \(g\in \cR\).
    For each \(k\) let \(f_k\) be such that \(g_k = (\id -A)f_k\).
    This means that \(f_k =  g_k + A f_k\) and, iterating, that \(f_k = \sum_{j=0}^{n-1} A^{j} g_k +  A^{n} f_k\) for any \(n\in \bN\).
    Observe that the set \(\cG:={\{g_k\}}_{k} \cup \{g\}\) is compact and consequently that \(\sum_{j=0}^{n-1} A^{j} \cG\) is compact.
    Moreover, since \( \lim_{n}{\nc{A^n}}^{{1}/{n}} < 1\), for any \(\epsilon>0\), there exists \(n\in \bN\) such that \(A^{n} \left( \cup_{k}f_k \right)\) can be covered by a finite number of \(\epsilon\)-balls.
    This means that there exists a subsequence of \(f_{k}\) which converges to some \(f\) and so \(g = (\id -A)f \in \cR\).

    Finally we show that \uline{\(\Coker(\lambda \id -A )\) is finite dimensional}.
    Since \(\Ran(\id -A )\) is closed, an application of the Hahn-Banach theorem implies that its codimension equals the dimension of \(\Ker(\id -A^{*})\) in \(\cB^{*}\).
    By Lemma~\ref{lem:dual-nc} the first paragraph applies.
\end{proof}


The essential spectral radius is defined as 
\[
    \EssRad(A) := \sup \left\{ \abs{z} : z \in \EssSpec(A) \right\}.
\]

\begin{theorem}
    Let \(A\) be a bounded operator on \(\cB\).
    Then  
    \[
        \EssRad(A) = \lim_{n\to \infty}{\nc{A^n}}^{{1}/{n}}.
    \]
\end{theorem}
\begin{proof}
    Lemma~\ref{lem:is-fredholm} implies that \(\EssRad(A) \leq \lim_{n}{\nc{A^n}}^{{1}/{n}}\).
    For the other part....
\end{proof}

%%%%%%%%%%%%%%%%%%%%%%%%%%%%%%%%%%%%%%%%%%%%%%%%%%%%%%%%%%%%%%%%%%%%%%%%%%%%%

\section*{Notes}

\begin{enumerate}
    \item The notation and definitions used here match Davies' book~\cite{Davies07}.
    \item Lemma~\ref{lem:is-fredholm} is a generalization of \cite[Lemma 4.3.3]{Davies07}.
    \item Lemma~\ref{lem:dual-nc} is \cite[Lemma 2]{Nussbaum70}.
    \item The spectral decomposition of an operator using spectral projectors and a remainder operator with spectral radius arbitrarily close to the essential spectral radius is implied by \cite[Theorem 4.3.18]{Davies07}.
    \item A good overview of measures of non-compactness and the essential spectrum can be found in \cite[\S I.4]{EE87}.
\end{enumerate}

%%%%%%%%%%%%%%%%%%%%%%%%%%%%%%%%%%%%%%%%%%%%%%%%%%%%%%%%%%%%%%%%%%%%%%%%%%%%%


\begin{thebibliography}{10}
    % spell-checker: disable 

    \bibitem{Davies07}
    E.B. Davies.
    \newblock Linear Operators and their Spectra.
    \newblock \emph{Cambridge studies in advanced mathematics}, 106. Cambridge University Press (2007).

    \bibitem{EE87}
    D.E.~Edmunds \& W.D.~Evans. 
    Spectral Theory and Differential Operators. 
    \emph{Oxford Clarendon Press} (1987).

    \bibitem{Hennion93}
    H.~Hennion.
    \newblock Sur un th\'{e}orem spectral et son application aux noyaux Lipchitziens,
    \newblock \emph{Proc. Amer. Math. Soc.} 118 (1993) 627--634.

    \bibitem{Nussbaum70}
    R.D.~Nussbaum.
    \newblock The radius of the essential spectrum.
    \newblock \emph{Duke Math. J.} 37 (1970) 473--478.

    % spell-checker: enable 
\end{thebibliography}

\end{document}
